\documentclass[12pt]{article}
\usepackage[utf8]{inputenc}
\usepackage[T1]{fontenc}
\usepackage{hyperref}
\title{Mayaguduna Language Resources for K-12 Curriculum Development}
\author{Compiled by Claude Sonnet 5.5 (claude-sonnet-5-5) for Andr\'as Kornai}
\date{}
\begin{document}
\maketitle

\section{Introduction}
Mayaguduna (ISO 639-3 xmy) is also spelled Mayi-Kutuna or Maikudunu. It is a Mayi (Mayabic) language of the Gulf country of northern Queensland. Wikipedia calls it extinct, while the Endangered Languages Project rates it critically endangered. I found no evidence of living speakers or of any school programme. The search found only a few sources, so this report is short.

\section{Primary Educational Resources}
I found no primers, textbooks or school materials for this language.

\section{Linguistic and Reference Materials}
\begin{itemize}
\item Wikipedia, Mayi-Kutuna language: \url{https://en.wikipedia.org/wiki/Mayi-Kutuna_language}
\item Wikipedia, Mayi-Kutuna people: \url{https://en.wikipedia.org/wiki/Mayi-Kutuna_people}
\item Glottolog reference entry for Gavan Breen (1981), the main documentation of the Mayi languages: \url{https://glottolog.org/resource/reference/id/64840}
\item Glottolog reference entry for Sharpe (1983): \url{https://glottolog.org/resource/reference/id/2586}
\item Endangered Languages Project entry: \url{https://endangeredlanguages.com/lang/3237?hl=de}
\end{itemize}

\section{Religious and Community Resources}
I found no scripture translations or community organisations specific to this language.

\section{Digital and Audio Resources}
\begin{itemize}
\item PARADISEC catalogue, Gavan Breen collection GB30, item 001. The search summary describes it as 1970s vocabulary and elicitation recordings made with speakers in Queensland: \url{https://catalog.paradisec.org.au/collections/GB30/items/001}
\end{itemize}

\section{Curriculum Development Context}
Little material exists, and it is archival and linguistic. A heritage or revitalization curriculum would have to be built from Breen's published work and the PARADISEC recordings. It should be developed with descendant communities and the relevant archives. Access conditions for the recordings were not checked.
\end{document}